% A145: a uniform finite-local obstruction and an injective binary construction.
\subsection{One finite quotient detecting a continuum of completed pairs}

\begin{statement}{A.145}{Proved}
{A uniform nine-place obstruction and a continuum of ordinary completed pairs}
Retain the actual curve and quotient of A.139--A.141,
\[
 \begin{gathered}
 D^\circ:\quad aW^4=T(T-1)(T-\tfrac12),\qquad
 t=4T(1-T),\\
 a=\frac{318039202706815117780707}{27871484},\qquad
 \overline Q_D=\Pi/\overline K_D
       =\bigl(\F_2[\tau]/\tau^6\bigr)\rtimes C_8.
 \end{gathered}
\tag{A.145.1}
\]
Put
\[
 S=\{2,3,7,13,17,23,139,191,1119169\},
\tag{A.145.2}
\]
and prescribe, only at these primes, the ordinary tripod centres
\(x_p=4T_p(1-T_p)\) from (A.134.11).
There is no continuous section of
\(E_U/\overline K_D\to G_\Q\) whose combined Kummer class is
\[
              2\kappa(t)+\kappa(1-t)=[64a^2]
                       \quad\text{in }H^1(G_\Q,\mu_8),
\tag{A.145.3}
\]
whose local image at each \(p\in S\) is geometrically conjugate
to the image of the prescribed ordinary point section, and whose
image at every other place is geometrically conjugate to an
ordinary or cuspidal local section.  No centre outside \(S\)
is prescribed in this exclusion.

There is an injection
\[
 \{0,1\}^{\mathbf N}\longrightarrow
        (\widehat{\Q^\times})^2,
        \qquad \epsilon\longmapsto(A_\epsilon,B_\epsilon),
\tag{A.145.4}
\]
with the following properties.  Every image is an ordinary
completed unit-equation pair, both coordinates have infinitely
many nonzero ordinary integer valuations, its local centres at
\(S\) are the prescribed ones, and
\([A_\epsilon^2B_\epsilon]_8=[64a^2]_8\).
Every image is excluded by the same quotient (A.145.1), of
order \(512\) and exact class six, with these local images.
In particular no full Selmer section has either image pair.
The actual projections from the full metabelian and class-six
quotients in (A.141.3) give the same locally matching exclusion.

For all these pairs the extra coordinate at
\(b=-715208/36481\) also has the stronger property of A.142:
its ordinary local values are units outside the one fixed set
\(S_0\) of (A.142.1), but their quadratic reciprocity sum
against \(d=18512174429\) is \(1/2\).
There are exactly continuum many pairs in the image.
With the product topology on the binary sequences and the
inverse-limit topology of discrete power quotients on the
completed coordinates, this injection is continuous and is
a homeomorphism onto its compact image.
This statement proves a uniform detector for the displayed
family; it makes no claim that this quotient detects every
infinite-support pair.
\end{statement}

\begin{proof}
\emph{The proper class at every outside-set cusp.}
The proper class of A.139 is
\[
               \alpha_D=(T-c,d),\qquad c=-338/191.
\tag{A.145.5}
\]
The Chern-class correction and exact central-extension map
(A.141.8)--(A.141.13) give a class
\[
 e_D\in H^2(E_D/\overline K_D,\mu_2),\qquad
 E_D=\pi_1(D^\circ),
\tag{A.145.6}
\]
whose pullback at a proper point is its \(\alpha_D\)-value.
This also holds for any local cusp packet after its map to the
proper boundary point, because that map kills boundary inertia
and the correcting line bundle has zero Chern class on a field.

For completeness the four boundary values are literal local
quaternions.  At \(T=0,1,h=1/2\) they are
\[
 (338/191,d),\qquad (529/191,d),\qquad(867/382,d).
\tag{A.145.7}
\]
At infinity use the local parameter \(w=W/T\).
The equation in (A.145.1) gives
\[
 w^4(T-c)=a^{-1}(1-1/T)(1-h/T)(1-c/T).
\tag{A.145.8}
\]
The right side is a unit at infinity with residue \(a^{-1}\).
Multiplication by \(w^4\) does not change the quaternion class,
so the boundary value is \((a^{-1},d)\).
All four first entries and \(d\) are units at every \(p\notin S\).
Such \(p\) is odd, and the quadratic character of \(d\) is
unramified; its reciprocity pairing with a unit is zero.
Thus every one of the four boundary values is zero outside \(S\).
A.139 proves zero at every ordinary \(D^\circ(\Q_p)\)-point
outside this same set, by its exhaustive valuation calculation.
At the real place all values are zero since \(d>0\).

\emph{The actual local comparison for the finite quotient.}
Suppose a section with the displayed data exists.
The retained combined character in (A.141.14)--(A.141.18),
with (A.145.3), gives a fixed sheet of the actual degree-eight
cover.  After geometric conjugation the section therefore has
image in \(E_D/\overline K_D\).
Pulling back (A.145.6) gives a class in
\(H^2(G_\Q,\mu_2)=\operatorname{Br}(\Q)[2]\).

We verify its local interpretation even at a cuspidal place.
Choose a representative local ordinary or cusp section whose
image in \(E_U/\overline K_D\) is conjugate to the given one,
and lift its finite geometric conjugator to \(\Pi\).
After this conjugation its full image lies in \(E_D\): the
inverse image of \(E_D/\overline K_D\) in \(E_U\) is exactly
\(E_D\), since \(\overline K_D\subset H_D\).
For an ordinary section the fixed sheet is an actual
\(\Q_p\)-point above its centre.  For a cusp section its fixed
sheet belongs to an inertia orbit stable under \(G_{\Q_p}\).
Boundary points of the normalized cover are precisely these
inertia orbits; the stable orbit is therefore a rational
boundary point.  The intersection of its original cusp
decomposition group with \(E_D\) is that boundary decomposition
group.  Its map to \(\pi_1(\overline D)\) factors through the
residue Galois group, with inertia killed.  This proves the
claimed evaluation through (A.145.6) without a restriction on
the tangential parameter.

At the primes of \(S\), A.139 computes every possible local
sheet above the prescribed centre.  The value is \(1/2\) at
two and zero at the other eight primes.  At every remaining
finite place the preceding ordinary and boundary calculations
give zero, and the real value is zero.  The sum is \(1/2\),
contrary to global reciprocity.  This proves the first assertion
and its consequences through the actual quotient maps of A.141.

\emph{The initial state for the binary construction.}
We use the initial state of A.142, including every prescribed
point at its finite set \(S_0\), its real point over
\(R=529/382\), and its modified point at \(p_*=4679\).
At each later state the actual mixed curve is
\[
       C_{\boldsymbol n}:\quad T-i=a_iX_i^{n_i},
                   \qquad i=0,1,h,\qquad 4\mid n_i.
\tag{A.145.9}
\]
The three fixed common boundary fields are exactly those of
(A.142.11).  The controlled set is finite and contains all
coefficient, exponent, boundary-difference and ramification
primes.  We retain every controlled centre and its actual
roots.  Both \(T-c\) and \(T-c'\), where \(c'=1-c\),
are units at all controlled primes outside the initial \(S_0\).
Outside the current controlled set there exist unit-root local
points with these two differences units.  This last assertion
is the exact mixed-cover bound (A.142.12)--(A.142.13).
The first signed refinement uses the root-symbol sums in
(A.134.21); after that refinement each root-symbol value is zero
individually, as in A.142.

The nodes are indexed by finite binary words.  At depth \(j\)
we have completed the preceding cycles of factors
\(2,3,\ldots,j+1\); the next factor is \(m=j+2\).
Before refining coordinate zero, enlarge the controlled set
as in (A.90.20)--(A.90.26), using the bound for
\((mn_0,n_1,n_h)\).  Retain at newly controlled primes the
unit-root points with both additional differences units.
There are still only finitely many controlled primes.

Let \(L\) be the splitting field of the current boundary
root algebra at \(T=0\):
\[
 \Q[Z_1^{\pm1},Z_h^{\pm1}]/
       (Z_1^{n_1}+1/a_1,\ Z_h^{n_h}+h/a_h).
\tag{A.145.10}
\]
Choose a fresh prime \(q\) exceeding the indicated new bound,
outside the controlled set and all nonzero differences among
\(0,1,h,c,c'\), and completely split in the finite Galois
compositum of \(L\) and the three fields (A.142.11).
Chebotarev, in the exact identity-conjugacy-class form of
\cite[Theorem 8.31]{AgentMilneANT}, supplies infinitely many
such primes.  The same chosen \(q\) is used for the two
children of this node.

\emph{The two actual local choices.}
For child \(1\) prescribe
\[
                        T_q=a_0q^{n_0},\qquad X_0=q.
\tag{A.145.11}
\]
The other two coordinates have unit residue roots at \(T=0\)
by (A.145.10).  Their derivatives are units, so they have
actual Hensel lifts at the displayed \(T_q\).
All coefficients are \(q\)-units.  Therefore
\[
     v_q(T_q)=n_0>0,\qquad v_q(T_q-1)=0,
       \qquad T_q-c,T_q-c'\in\Z_q^\times.
\tag{A.145.12}
\]
Every root symbol against its fixed boundary field is zero:
the prime splits each such field, including the field for
the nonunit root \(X_0=q\).

For child \(0\), the point bound for the current curve with
both extra fibres removed gives a residue point with all
root coordinates and \(T-c,T-c'\) units.  Choose a rational
lift of its centre in that residue class and lift the roots
by their unit derivatives.  This is an actual ordinary point
with
\[
          v_q(T_q)=v_q(T_q-1)=0,
                 \qquad T_q-c,T_q-c'\in\Z_q^\times.
\tag{A.145.13}
\]
Its root symbols also vanish.  In either child add \(q\) to
the controlled set, retaining its newly chosen point.

\emph{The nonunit-root refinement and preservation of the bit.}
We spell out why the nonunit root in (A.145.11) is permitted.
Let \(S'\) denote the enlarged controlled set and let
\(x_p=X_0(P_p)\) be its exact coordinate-zero targets.
Set
\[
          d_0=\varepsilon_\infty
                    \prod_{p\in S'}p^{v_p(x_p)},
       \qquad \varepsilon_\infty=
       \begin{cases}-1,&\text{first signed step},\\
                     1,&\text{later steps}.
       \end{cases}
\tag{A.145.14}
\]
Each \(x_p/d_0\) is a unit.  Its sufficiently precise unit
CRT class acts trivially on the fixed boundary common field.
Indeed subtract the global reciprocity identity for the
rational number \(d_0\) from the zero total root-symbol sum.
Outside \(S'\) the number \(d_0\) is a unit in every
unramified boundary field; its real sign matches the prescribed
root.  The resulting unit sum is zero.  The local cyclotomic
reciprocity map is the inverse-unit map, so this is precisely
the trivial restriction of the CRT class, as in
(A.134.23)--(A.134.24) and (A.90.20)--(A.90.23).

The actual stabilizer-image equality (A.90.25) writes this CRT
class as a product of images of boundary-root stabilizers.
Choose the associated distinct Chebotarev primes outside
\(S'\), and also the extra completely split prime required
in (A.90.26).  Their product times \(d_0\) is the rational
multiplier \(u\).  For sufficiently large precisions,
\(x_p/u\) has a principal \(m\)-th root \(Y_0\) in
\(\Q_p\) for every \(p\in S'\).  Thus the exact map is
\[
 \begin{gathered}
 X_0=uY_0^m,\qquad n'_0=mn_0,
                \qquad a'_0=a_0u^{n_0},\\
       T=T,
       \qquad X_1=Y_1,\quad X_h=Y_h.
 \end{gathered}
\tag{A.145.15}
\]
All three boundary fields remain unchanged because the
coefficient changes are fourth powers, with the exact
mixed-exponent comparison of A.90.  The new controlled root
\(Y_0\) is a principal local square and hence has zero symbols.
At the real place it is positive, including after the first
negative multiplier.  At a new coefficient prime the explicit
centre near \(T=0\) in A.142 has both additional differences
units and all new roots units.  At all other primes the
mixed-cover point bound provides the required points.

The new prime \(q\) is controlled before choosing the
Chebotarev primes.  Consequently
\[
     v_q(u)=v_q(x_q)=
          \begin{cases}1,&\text{child }1,\\0,&\text{child }0,
          \end{cases}
        \qquad v_q(Y_0)=0.
\tag{A.145.16}
\]
The old centre at \(q\) is retained literally.  All subsequent
refinements at that controlled prime have unit target roots,
and so have unit multipliers there.  Equations (A.145.12)
and (A.145.13) are therefore permanent along the respective
branches.  Finish this cycle by refining coordinates one and
\(h\) with the same factor \(m\), using the identical
controlled-place rule.  Keep one extra completely split prime
for each coordinate refinement, in addition to the binary
prime, exactly as in A.142.  This ensures infinite support even
along the all-zero binary path.

\emph{The completed maps and injectivity.}
Fix the countable choices at every node: order rational
centres, residue roots, admissible primes and finite
stabilizer decompositions, and take the first allowed choices.
Hensel branches are then uniquely determined by their specified
residue and precision conditions.  The preceding existence
proof makes these choices at every finite node.  Every infinite
binary path has common exponents \(4m!\) after its factor
\(m\) cycle.  The coefficient identity in (A.145.15) therefore
gives compatible completed coordinates
\(A_0,A_1,A_h\), with their fixed initial classes modulo four.
Every rational prime is eventually controlled and retains its
actual centre.  Passing to any fixed finite exponent in the
late equations of (A.145.9) gives, literally,
\[
 \begin{gathered}
       (A_i)_p=\kappa_p(T_p-i),\\
       A=-4A_0A_1,\qquad B=4A_h^2,\qquad x_p=4T_p(1-T_p).
 \end{gathered}
\tag{A.145.17}
\]
These are continuous maps in the exponent multiplicative
completions, and the displayed identities prove the ordinary
unit equation at every finite place.  The real signs remain
those of the ordinary centre \(4R(1-R)\).
The initial coefficient product is exactly \(a=a_0a_1a_h\),
and hence
\[
        [A^2B]_8=[64(A_0A_1A_h)^2]_8=[64a^2]_8.
\tag{A.145.18}
\]
Here squaring gives the explicit homomorphism
\(\Q^\times/(\Q^\times)^4\to
  \Q^\times/(\Q^\times)^8\) used in the last equality.

Two binary paths with first differing bit at a node share
its prime \(q\).  At that prime (A.145.17) and the permanent
centre values give
\[
             v_q(A)=n_0\quad\text{on child }1,
               \qquad v_q(A)=0\quad\text{on child }0.
\tag{A.145.19}
\]
The valuation map here is the inverse limit of the ordinary
valuation maps modulo every exponent, with values in
\(\widehat\Z\).  The two displayed integers are distinct
there.  This proves injectivity of (A.145.4).
The extra coordinate-zero and coordinate-one primes have
respectively positive valuations of \(T\) and \(T-1\),
with the other factor a unit, and thus give infinitely many
nonzero valuations of \(A\).  The extra coordinate-\(h\)
primes give positive valuations of \(B\), by its square in
(A.145.17).  All are retained by the same unit-multiplier rule.
Since each \(\Q^\times/(\Q^\times)^n\) is countable, its
countable inverse limit has cardinality at most continuum.
The injection supplies equality for the displayed family.
For a fixed exponent \(N\), choose a cycle with \(N\mid4m!\).
All extensions of a binary prefix through that cycle have
the same three coefficient classes modulo \(N\), by
(A.145.15).  Their pairs consequently have the same image in
\((\Q^\times/(\Q^\times)^N)^2\).  This proves continuity
of the binary map into the inverse limit of these discrete
quotients.  The binary sequence space is compact, and the
target is Hausdorff.  A continuous injection between these
spaces is a homeomorphism onto its compact image: images of
closed subsets are compact and hence closed in the target.
This proves the asserted topology of the constructed family.

Every initial point of A.142 is retained, and every point
introduced outside its \(S_0\) has both differences units.
The literal identity
\[
                x_p-b=-4(T_p-c)(T_p-c')
\tag{A.145.20}
\]
therefore preserves its complete invariant table: \(1/2\)
at \(4679\) and zero elsewhere.  It also preserves the
units outside \(S_0\) and the real value.  Thus each family
member has the additional-coordinate obstruction asserted.

Finally, a full Selmer section with one of these pairs would
have the specified ordinary centre at every finite prime.
The local Kummer map \(\Q_p^\times\to
\widehat{\Q_p^\times}\) is injective: valuation, the finite
unit group and the principal-unit group all inject into their
full power completions.  At cusp zero the coordinate \(B\)
is one, at cusp one \(A\) is one, and at infinity \(B/A=-1\).
For an ordinary pair \((x_p,1-x_p)\) these equalities would
respectively force \(x_p=0\), \(x_p=1\), or \(1=0\), by
this injectivity.  Each cusp alternative is therefore excluded.
An ordinary local section has unique centre \(x_p\) by the
same injection.  Its image in (A.145.1) would satisfy all the
already excluded local conditions.  This proves the full
Selmer assertion with the exact direction of the comparison.
\end{proof}

The first binary node has an executed finite witness
\[
 \begin{gathered}
 q=1065217>1065024,\qquad T^{(0)}=36,\\
 (X_0,X_1,X_h)^{(0)}\bmod q=(125292,99176,89494),\\
 T^{(1)}=\tfrac{529}{382}q^4,
       \qquad (X_1,X_h)^{(1)}\bmod q=(299139,97645).
 \end{gathered}
\tag{A.145.21}
\]
The independent program
\path{code/uniform_detector_branch_independent_audit.py}
checks primality, all displayed fourth-power equations, the
strict bound and splitting of the five boundary quadratic
characters.  The strengthened root certificate also checks
the two extra unit conditions, the literal valuation gap,
and one actual Hensel step for every residue root.
These are certificates for this finite node.  The proved
Chebotarev and refinement construction supplies the infinite
binary tree and its compatible limits.
